% !TEX root = ../../../main.tex
% Appendix entry: detection.dehshibi_spikes -- deprecated monolith (Appendix A).
\entryfield{Algorithm}
\code{_run} in \code{Adapters/detection_dehshibi_spikes.py} calls \code{detect_spikes} in
\code{Working/Detection/analysis/dehshibi_authors.py} with $\texttt{thr} = \texttt{epsilon\_factor}$
and a separation of $\texttt{n\_p}/f_s$ seconds. \code{detect_spikes} is the three template
stages in one call:
\begin{enumerate}
\item \code{transform}: the scaled Morse scalogram in 3000\,s windows (stage~1,
  \blockentry{preprocessing.wavelet_transform});
\item $\Omega$ = the sum of the rows whose frequency is at most a quarter of the frequency
  range, chosen by \code{low_band_rows} from the actual frequencies (stage~2,
  \blockentry{detection.wavelet_summation});
\item \code{detect_from_omega}: candidate regions, the spike / pseudo-spike split, envelope
  regions, merge (stage~3, \blockentry{detection.summation_threshold}).
\end{enumerate}
Since fixup-J the block and the template share this core, so they cannot drift apart. It
follows the authors' MATLAB, not the printed paper; the differences are listed under
stage~3. \emph{min\_spike\_duration} and \emph{min\_roi\_wavelet} belonged to the printed
Algorithm~4, have no counterpart in the authors' code, and are accepted and ignored.

The one real difference from the template is step~2: the monolith picks the band from the
frequencies, so it is the authors' band at any sampling rate, whereas the template passes a
fixed fraction (0.7586) that is the authors' band only for 3000-sample windows.

\entryfield{Assumptions}
\begin{itemize}
\item \emph{Units.} As for the template: the spans do not depend on the unit (every threshold
  is relative).
\item \emph{Sampling rate.} \emph{n\_p} is in samples (the template's
  \emph{min\_separation\_s} is in seconds); 60 is the authors' value at 1\,Hz. The window is
  fixed at 3000\,s and state slicing is off; neither can be changed here.
\item \emph{Missing values.} As for the template: a NaN blanks its whole 3000\,s window, which
  is then skipped.
\end{itemize}

\entryfield{Outputs}
A \texttt{SpanSet} of half-open spans, \emph{without} labels or scores. \code{meta} holds
\code{pseudo_spikes} (half-open), \code{n_spikes}, \code{n_pseudo_spikes} and \code{n_chunks};
none of the funnel the template's last block reports. A \code{_plot} hook draws spikes and
pseudo-spikes with \code{plot_spike_detection} from the older module
\code{dehshibi_detection_analysis.py}, with $f_s$ fixed at 1.

\entryfield{Cost}
As the template's stages together: FFT-bound, about $n\log n$ (cost model exponent 1.1; the
estimate is empty until \code{Working.block_cost.calibrate()} has run). The block declares no
\code{max_span_samples}, so, unlike \blockref{preprocessing.wavelet_transform}, a long span is
not refused: the float32 image costs 4 bytes per row per sample, \unverified{about 0.9\,GB for
2.6 million samples at 87 rows}.

\entryfield{Failure modes}
\begin{itemize}
\item Every failure mode of the three stages, above all the lack of selectivity on real
  recordings (\blockentry{detection.summation_threshold}).
\item \emph{No memory guard} (above).
\item \emph{Settings that do nothing.} Two of its four settings are ignored, with no warning
  beyond their descriptions.
\end{itemize}

\entryfield{Validation}
\code{tests/test_template_dehshibi.py} runs the template and this block on the 3000\,s, 1\,Hz
fixture recording and checks that both return exactly the five spikes the authors' MATLAB
returned. The core is checked stage by stage in \code{tests/test_dehshibi_authors.py}. The
detection rates on synthetic and real data are given under
\blockentry{detection.summation_threshold}; they apply here unchanged.

\entryfield{References}
\citet{dehshibi2021electrical}; the authors' code, Zenodo doi:10.5281/zenodo.3997031;
\code{docs/prompts/fixup/reports/J-dehshibi-vs-the-paper.md}.

\entryfield{Tests}
\code{tests/test_template_dehshibi.py} (agreement with the template and with MATLAB; filed under
control and its description points to the template); \code{tests/test_t06_adapter_remap.py}
(the output is a well-formed span set).
